\subsection{Rancangan Skenario Pengujian}
\label{subsec:rancangan-skenario-pengujian}

Rancangan skenario pengujian diturunkan dari rancangan alur akses dan delegasi pada \texttt{R-TA-08}. Fokus pengujian adalah membuktikan bahwa setiap aktor hanya memperoleh hak minimum yang diperlukan untuk menjalankan tugasnya, serta bahwa token turunan tidak dapat memperluas cakupan token induk.

Skenario uji dirancang dalam bentuk alur positif dan negatif. Alur positif memastikan aktor dapat menjalankan operasi yang memang termasuk dalam cakupan token. Alur negatif memastikan \textit{Server} PRE menolak permintaan ketika aktor mencoba membaca, menulis, atau mendelegasikan data di luar \textit{caveats} yang diberikan.

\begin{table}[!ht]
	\centering
	\caption{Rancangan Skenario Uji Berbasis Aktor}
	\label{tab:rancangan-skenario-uji-aktor}
	\small
	\begin{tabular}{|p{0.12\textwidth}|p{0.48\textwidth}|p{0.23\textwidth}|}
		\hline
		\textbf{ID}       & \textbf{Skenario}                                                         & \textbf{Target Rancangan}                            \\ \hline
		\texttt{PS-TA-01} & Pasien memberikan token akses awal kepada petugas rekam medis.            & \texttt{R-TA-03}, \texttt{R-TA-08}                   \\ \hline
		\texttt{PS-TA-02} & Petugas rekam medis mendelegasikan akses pemeriksaan awal kepada perawat. & \texttt{R-TA-04}, \texttt{R-TA-08}                   \\ \hline
		\texttt{PS-TA-03} & Perawat menulis segmen anamnesis atau pemeriksaan awal.                   & \texttt{R-TA-02}, \texttt{R-TA-05}, \texttt{R-TA-06} \\ \hline
		\texttt{PS-TA-04} & Perawat mencoba menulis terapi.                                           & \texttt{R-TA-05}, \texttt{R-TA-06}                   \\ \hline
		\texttt{PS-TA-05} & Dokter membaca dan menulis diagnosis serta terapi.                        & \texttt{R-TA-02}, \texttt{R-TA-05}, \texttt{R-TA-06} \\ \hline
		\texttt{PS-TA-06} & Dokter mendelegasikan akses terbatas kepada laboratorium.                 & \texttt{R-TA-04}, \texttt{R-TA-08}                   \\ \hline
		\texttt{PS-TA-07} & Petugas laboratorium membaca permintaan dan menulis hasil pemeriksaan.    & \texttt{R-TA-02}, \texttt{R-TA-05}, \texttt{R-TA-06} \\ \hline
		\texttt{PS-TA-08} & Petugas laboratorium mencoba membaca diagnosis.                           & \texttt{R-TA-05}, \texttt{R-TA-06}                   \\ \hline
		\texttt{PS-TA-09} & Dokter mendelegasikan akses terbatas kepada apoteker.                     & \texttt{R-TA-04}, \texttt{R-TA-08}                   \\ \hline
		\texttt{PS-TA-10} & Apoteker membaca resep dan menulis dispensing.                            & \texttt{R-TA-02}, \texttt{R-TA-05}, \texttt{R-TA-06} \\ \hline
		\texttt{PS-TA-11} & Apoteker mencoba membaca anamnesis atau hasil laboratorium.               & \texttt{R-TA-05}, \texttt{R-TA-06}                   \\ \hline
		\texttt{PS-TA-12} & Token kedaluwarsa digunakan.                                              & \texttt{R-TA-03}, \texttt{R-TA-05}                   \\ \hline
		\texttt{PS-TA-13} & Token digunakan oleh wallet yang tidak sesuai.                            & \texttt{R-TA-04}, \texttt{R-TA-05}                   \\ \hline
		\texttt{PS-TA-14} & Delegasi melebihi \texttt{max\_delegation\_depth}.                        & \texttt{R-TA-04}, \texttt{R-TA-05}                   \\ \hline
		\texttt{PS-TA-15} & Token yang telah dicabut digunakan kembali.                               & \texttt{R-TA-07}                                     \\ \hline
	\end{tabular}
\end{table}

\texttt{R-TA-01} bersifat lintas-komponen sebagai rancangan arsitektur umum.
Oleh karena itu, rancangan ini tidak diuji melalui satu skenario khusus,
melainkan dinilai melalui keberhasilan implementasi \texttt{I-TA-01} sampai
\texttt{I-TA-06} dan keseluruhan pengujian sistem.

Rancangan tersebut menjadi acuan bagi pengujian implementasi pada Bab \ref{bab:implementasi-dan-evaluasi}. Skenario dengan ekspektasi diterima digunakan untuk memverifikasi kelengkapan alur layanan, sedangkan skenario dengan ekspektasi ditolak digunakan untuk memverifikasi prinsip \textit{least privilege}.